\documentclass[a4paper,12pt]{article}

% Pacchetti comuni
\usepackage[utf8]{inputenc}    % Codifica UTF-8
\usepackage{amsmath}            % Simboli matematici
\usepackage{amsfonts}           % Font matematici
\usepackage{amssymb}            % Simboli matematici
\usepackage{graphicx}           % Gestione delle immagini
\usepackage{hyperref}           % Collegamenti ipertestuali

\title{Appunti della lez}
\author{Nicolo' Giuliani}
\date{\today}

\begin{document}
\section{Ripasso}
Il determinante e' una funzione $det \text{ : }M_{nn}(\mathbb{R}) \to \mathbb{R}$ che gode di 4 proprieta'.
\\ovvero $a,b,c$ e che il det A = det $A^T$ (\textit{Transposta di A}).
\\Nota: le proprieta vengono sostituendo la parola righa con la parola colonna.
\subsection{Teorema di Binet}
$det AB=detAdetB$
\subsection{Calcolo del determinante}
1) Gauss
\\2) Laplace
\\\subsection{Matrice inversa}
Def: $A \in M_{mn}(\mathbb{R})$ si dice invertibile se esiste $B \in M_{nn}(\mathbb{R})$ tale che $AB=BA=I$.
\\Proposizione: $A \in M_{nn}(\mathbb{R})$  e' invertibile $\iff det A \neq 0$
\\Proposizione: Sia $A \in M_{nn}(\mathbb{R})$  e sia $B \in M_{nn}(\mathbb{R})$  tale che $AB=I$. Allora $B$ e' l'inversa di $A$, cioe' $BA=I$.
\\Dim: $AB=I$ $= det(AB)=det(I)=1$ attraverso binet $\to det(A)det(B)=1 \to det(A)\neq 0$, quindi A e' invertibile e $A^{-1}$ l'inversa. 
\\Moltiplichiamo entrambi i membri di $AB=I$ a sinistra per $A^{-1}$. Quindi $A^{-1}AB=A^{-1}I \to IB=A^{-1} \to B=A^{-1}$ in particolare $BA=I$.
Proposizione (versione bis): $BA=I \to B=A^{-1}$
\section{Calcolo dell'inversa}
(Osservazione: $det(A) \neq 0$) \\
1) Laplace: $b_{ij}=\frac{1}{detA}T_{ij}$
\\2) Algoritmo di Gauss: \\
Si scrive $A \mid I$ 
\\ con le operazioni elementari di Gauss arriviamo a $I \mid B$ allora $B=A^{-1}$
\\Esempio:
\begin{align*}
    A=\begin{pmatrix}
        1&2&0\\-1&3&1\\0&1&-2
    \end{pmatrix} \\
    A\mid I=\begin{pmatrix}
    1 & 2 & 0 & | & 1 & 0 & 0 \\
    -1 & 3 & 1 & | & 0 & 1 & 0 \\
    0 & 1 & -2 & | & 0 & 0 & 1
    \end{pmatrix} \\
    \text{la portimo a scala} \\
    \begin{pmatrix}
        1&2&0&|&1&0&0\\
        0&1&-2&|&0&0&1\\
        0&0&11&|&1&1&-5
    \end{pmatrix}
\end{align*}
Quando la matrice e' a scala faccia comparire tutti gli "1" sulla diagonale principale.
\begin{align*}
    \begin{pmatrix}
        1&2&0&|&1&0&0\\
        0&1&-2&|&0&0&1\\
        0&0&1&|&\frac{1}{11}&\frac{1}{11}&\frac{-5}{11}
    \end{pmatrix}
\end{align*}
Ripartiamo con Gauss dal basso utilizzando il pivot in posizione "n,n" (in questo caso 3,3), per far apparire tutti gli "0" sopra di esso.
\begin{align*}
    (R_2+2R_3)
    \begin{pmatrix}
        1&2&0&|&1&0&0\\
        0&1&0&|&\frac{2}{11}&\frac{2}{11}&\frac{1}{11}\\
        0&0&1&|&\frac{1}{11}&\frac{1}{11}&\frac{-5}{11}
    \end{pmatrix}
\end{align*}
Poi si utilizza il pivot "1" in posizione "n-1,n-1" e cosi' via.
\begin{align*}
    (R_1-2R_2)
    \begin{pmatrix}
        1&0&0&|&\frac{7}{11}&\frac{-4}{11}&\frac{-2}{11}\\
        0&1&0&|&\frac{2}{11}&\frac{2}{11}&\frac{1}{11}\\
        0&0&1&|&\frac{1}{11}&\frac{1}{11}&\frac{-5}{11}
    \end{pmatrix}=I\mid B 
    \\\to B=A^{-1}
\end{align*}
RICORDA:nella matrici 2x2 il determinante ha un significato geometrico, ovvero l'area.
\\Se il det=0 vuol dire la matrice e' una trasformazione in uno spazio inferiore.
\pagebreak
Applicazione lineare $F : \mathbb{R}^n \to \mathbb{R}^n$, sia $F=L_A$. Supponiamo che $F$ sia isomorfismo (iniettiva e suriettiva), sia $G=F^{-1}$.
$G : \mathbb{R}^n \to \mathbb{R}^n $. \\$FoG=$identita' $\mathbb{R}^n$. $GoF=$identita' $\mathbb{R}^n$.
\begin{align*}
    \mathbb{R}^n (F)\to \mathbb{R}^n (G=F^{-1})\to \mathbb{R}^n 
\end{align*}
Sia B la matrice associato a G, $G=L_A$. Osserviamo che la matrice associata ad $id_{\mathbb{R}^n}$ e':
\begin{align*}
    \begin{pmatrix}
        1&0&0&0\\0&1&0&0\\\vdots\\0&0&0&1
    \end{pmatrix}
\end{align*}
$L_I=id$
\begin{align*}
    e_1 \to e_1 \\
    e_2 \to e_2 \\
    e_n \to e_n
\end{align*}
\begin{align*}
    \mathbb{R}^n (F)(A)\to \mathbb{R}^n (G=F^{-1})(B)\to \mathbb{R}^n 
    \\GoF=id \iff BA=I
\end{align*}
$BA=I \to B$ e' l'inversa di A.
\\Quindi alla funzione inversa di $F^{-1}$ di F e' associata la matrice inversa di $A^{-1}$ di A e viceversa.
\section{Teorema 7.6.1 (teoremone)}
Sia $F: \mathbb{R}^n \to \mathbb{R}^n$ applicazione lineare e sia A la matrice associata ad F rispetto alla base canonica in dominio e codominio.  (Cioe' $F=L_A$) 
\\Sono equivalenti queste condizioni:
\\1) F e' un insomorfismo
\\2) F e' iniettiva
\\3) F e' suriettiva
\\4) dim(Im F)=n
\\5) rk(A)=n
\\6) le colonne di A sono linearmente indipendenti
\\7) le righe di A sono linearmente indipendenti
\\8) il sistema $Ax=0$ ha un unica soluzione
\\9) $\forall b \in \mathbb{R}^n$ il sistema Ax=b ha un unica soluzione.
\\10) A e' invertibile
\\11) $det(A) \neq 0$
\\RICORDA: \textit{10 $\iff$ 11}
\\RICORDA: \textit{2 $\iff$ 1}
Dim: 
\begin{align*}
    F : \mathbb{R}^n \to \mathbb{R}^n \\
    F_{iniettiva}=Ker F = {0}  \to Ker F =0
    \\ \to dim(Im F)=dim(\mathbb{R}^n)-dim(Ker F)=n
    \\ \to F \text{ e' suriettiva} \to \text{ e' un isomorfismo}
\end{align*}
\\RICORDA: \textit{3 $\iff $ 1}
\begin{align*}
    F : \mathbb{R}^n \to \mathbb{R}^n \text{ e' suriettiva per ipotesi}\\ 
    \to dim(Im F)=n \to dim(Ker F)=dim \mathbb{R}^n
    \\ dim(Im F)=0 \to Ker F = {0} \to F \text{ e' iniettiva}
\end{align*}
Quindi $1 \to 3, 4 \to 2$ per def. di isomorfismo.
\\ $3 \iff 4$:
\begin{align*}
    F \text{ e' suriettiva} \iff Im F = \mathbb{R}^n
    \\ \iff (4.2.4) \text{  }dim(Im F)=dim(\mathbb{R}^n)
\end{align*}
$4 \iff 5$:
\begin{align*}
    Im F = \langle \text{colonne di A}\rangle \\
    dim(Im F)=n \iff dim(\text{rk(A)})=n
\end{align*}
$5 \iff 6$:
\begin{align*}
    rk(A)=n \iff dim(\langle \text{colonne di A}\rangle)=n 
    \\ \iff (GEL) \text{  le colonne di A sono linearmente indipendenti}
\end{align*}
$5 \iff 7$:
\begin{align*}
    rk(A)=n \iff dim(\langle \text{righe di A}\rangle) =n
    \\ \iff(GEL) \text{  le righe di A sono linearmente indipendenti}.
\end{align*}
$2 \iff 8$:
\begin{align*}
    F \text{ e' iniettiva} \iff Ker F={0}
    \\ \iff \text{ il sistema }Ax=0 \text{ ha unica soluzione. (quella nulla)}
\end{align*}
$1 \iff 9$:
\begin{align*}
    F \text{ e' suriettiva} \iff \forall b \in \mathbb{R}^n \mid F^{-1}(b) \text{ ha un solo elemento}
    \\ \iff \text{ il sistema Ax=b ha un unica soluzione}
\end{align*}
$1 \iff 10$:
\begin{align*}
    F \text{ e' isomorfiso } \iff \exists F^{-1} \\
    \iff \exists A^{-1} \iff A \text{ e' iniettiva}
    \\ \text{dimostrati prima del teoremone}
\end{align*}
ESERCIZIO:
\\stabilire per quali valori di k si ha che $F_K$ e' iniettiva
\begin{align*}
    F_K : \mathbb{R}^3 \to \mathbb{R}^3 \text{ e' appl. lin.}\\
    F_K(e_1)=e_1+ke_3 \\
    F_K(e_2)=ke_1+ke_2+6e_3 \\
    F_K(e_3)=e_1+ke_2+2ke_3
\end{align*}
\begin{align*}
    \begin{pmatrix}
        1&k&1\\
        0&k&k\\
        k&6&2k
    \end{pmatrix}
\end{align*}
2 modi per risolvere:
\\1)Sauss
\\2)Laplace
\\Faremo il secondo, scegliamo la prima colonna.
\begin{align*}
    det(A_K)=1(-1)^{1+1}det\begin{pmatrix}
        k&k\\6&2k
    \end{pmatrix}+0+(-1)^{3+1}kdet\begin{pmatrix}
        k&1\\k&k
    \end{pmatrix}
    \\ =2K^2-6K-K(K^2-K)
    \\ = K(K^2+K-6)
\end{align*}
\begin{align*}
    det(A)=0 \iff K=0 \text{ o } K^2+K-6=0 \to (K+3)(K-2) \to K=-3,K=2
\end{align*}

\end{document}